\documentclass[11pt]{article}
\usepackage[margin=25mm]{geometry}
\usepackage{fontspec}
\setmainfont{Latin Modern Roman}
\usepackage{amsmath,amssymb,mathtools}
\usepackage{unicode-math}
\setmathfont{Latin Modern Math}
\usepackage{xurl,hyperref}
\hypersetup{colorlinks=true,urlcolor=blue}
\setcounter{secnumdepth}{0}
\setlength{\parindent}{0pt}
\setlength{\parskip}{5pt}
\setlength{\emergencystretch}{3em}
\newcommand{\sref}[2]{\hyperlink{source:#1:#2}{[S#2]}}
\newcommand{\cataloguescope}[1]{\paragraph{Recorded target.}#1}
\newcommand{\R}{\mathbb R}
\newcommand{\C}{\mathbb C}
\newcommand{\N}{\mathbb N}
\newcommand{\Z}{\mathbb Z}
\newcommand{\Q}{\mathbb Q}
\newcommand{\F}{\mathbb F}
\newcommand{\eps}{\varepsilon}
\begin{document}
% UnsolvedMath ID 20002524; source catalogue code AIM-PHYSICS-0032
\setcounter{section}{20002523}
\section{Audible primitive two-torus areas on a flat three-torus}\label{20002524}
{\small UnsolvedMath ID 20002524 · Mathematical Physics}
\subsection{Definitions and mathematical statement}

\paragraph{Shared notation.}$\mathbb N=\{1,2,\ldots\}$; the standard number fields have their usual meanings. All additional conventions are stated locally.


\paragraph{Definitions and notation.}$\operatorname{Spec}(\Delta)$ is the spectrum of the Laplace operator for the chosen Riemannian metric and boundary condition. Minimal surfaces are critical points of area. An area spectrum requires specifying which minimal surfaces are included and how multiplicities or families are counted. In $T^4=T^2\times T^2$, $[T_1],[T_2]$ are the homology classes of the two torus factors; $g$ denotes the surface genus.

\paragraph{Catalogue wording retained for traceability.}
For three-manifolds, identify two-dimensional geometric objects related to $\operatorname{Spec}(\Delta)$ in a manner analogous to the geodesic-length relationship. Determine whether a suitable analogue of the minimal-surface area spectrum exists under appropriate geometric hypotheses. Compute the number $N(n_1,n_2,g)$ of genus-$g$ minimal surfaces in the flat $T^4=T^2\times T^2$ representing $n_1[T_1]+n_2[T_2]$.
\paragraph{Scope and source qualification.}The source leaves the counting convention and the relevant area-spectrum objects open; translations and continuous families cannot silently be discarded. The title’s primitive two-torus example does not replace the full three-part target. The extracted background remarks spill into separate subsequent Jardim Dirac-spectrum questions and the Hausel programme; those are not silently reassigned to this canonical three-part Cherkis excerpt. The catalogue title/status was paired with a machine-generated research attempt; that attempt is not peer-reviewed evidence of settlement. The statement covers the original source questions rather than only the title’s subcase.
\subsection{Short English statement}
Find surface-area analogues of spectral/length relations and count minimal surfaces in prescribed flat four-torus classes.
\subsection{Research attempt}
\paragraph{Model, date and investigated scope.}
OpenAI GPT-6 Astra (configured model \texttt{gpt-6-astra}), ultra reasoning; generated 11 October 2026. Derived primitive flat two-torus area versus Fourier eigenvalue and exhibited uncountably many translated genus-one minimal surfaces in a fixed T4 class. General area spectra and a well-defined N(n1,n2,g) remain unresolved.

\paragraph{Sources and literature check.}
Read the supplied complete definition, catalogue wording, scope qualifications and source list for this canonical ID. The original-source citations are inherited as supplied references; no fresh primary-source or novelty search was performed for this attempt. No current open-status claim is inferred from this calculation.

\paragraph{Flat three-torus area and eigenvalue relation.}Let \(T^3=\mathbb R^3/\Lambda\) have volume \(V\), and let \(k\) be a primitive vector in the dual lattice \(\Lambda^*\). The map \(x\mapsto k\cdot x\pmod1\) has connected flat two-torus fibers. Its gradient has constant length \(|k|\). Integrating over one period in the normal direction, whose physical length is \(1/|k|\), gives fiber area \(A_k=V|k|\). Each fiber is totally geodesic, hence minimal. The Fourier mode \(e^{2\pi ik\cdot x}\) has Laplace eigenvalue \(\lambda_k=4\pi^2|k|^2\), so
\[A_k=\frac{V}{2\pi}\sqrt{\lambda_k}.
\]
Primitivity matters: a nonprimitive \(lk\) defines a disconnected union of \(l\) parallel fibers and does not describe one primitive surface.
\paragraph{A counting obstruction in the four-torus clause.}In the product \(T^4=T^2\times T^2\), every \(T^2\times\{p\}\), with \(p\) in the second torus, is a distinct embedded flat minimal torus of genus one representing \([T_1]\). Thus the literal unquotiented count \(N(1,0,1)\) is uncountable. Translation equivalence, marked-point conditions or a virtual/moduli count would change the problem and must be declared, rather than silently imposed.
\paragraph{Gap.}The flat relation does not show that a general three-manifold Laplace spectrum determines minimal-surface areas, nor define a suitable area spectrum with multiplicities. The first missing step is a precise object/counting convention stable under the relevant geometry; for the four-torus question it must handle continuous families even in the simplest class. This elementary subcase and obstruction preserve all three source requests without asserting a universal spectral-surface theorem.

\paragraph{Verification and adversarial audit.}
Checked primitive dual-lattice fibers, normal period, eigenvalue factor 4 pi squared and the explicit genus-one translation family. This was an LLM adversarial audit, not independent human verification; no theorem prover was used. Prompts A then B applied to this ID: Checked primitive dual-lattice fibers, normal period, eigenvalue factor 4 pi squared and the explicit genus-one translation family. Scope remains limited: Specify area-spectrum and minimal-surface counting conventions that handle translation families, then prove a general spectral relation or enumerative formula.

Coordinator checked the primitive flat-torus fiber area and Fourier eigenvalue relation, together with the continuous family obstructing an unquotiented surface count. This additional check was model-assisted, not independent human verification.

\paragraph{Tools, process and contributions.}
Python 3 and local shell used to read supplied files and serialize this manuscript; mathematics derived symbolically in the text, with no numerical experiment or formal proof assistant. Contribution prompt A was followed by prompt B. GPT-6 Astra with ultra reasoning drafted the research and performed its stated self-audit; any additional coordinator contribution is identified in the verification and audit. A GPT-6-based Codex coordinator assembled the manuscript and checked repository integration; its exact variant and reasoning configuration were not exposed. No human mathematical contribution, independent human proof check or formal proof certification is supplied. Complete derivations and the consulted source references are included here; private reasoning transcripts are not published.

\paragraph{Final status and first remaining gap.}
\textbf{UNRESOLVED}. Specify area-spectrum and minimal-surface counting conventions that handle translation families, then prove a general spectral relation or enumerative formula.

\section{Completion Estimate}
\noindent\textbf{COMPLETION: 0\%}
The calculation is an elementary benchmark or conditional reduction; it establishes no claimed new portion of the complete source target.

\subsection{Sources}
\begin{itemize}
\item[\hypertarget{source:20002524:1}{S1}] ULAM.AI and the UnsolvedMath Contributors, UnsolvedMath, record 20002524 (AIM-PHYSICS-0032). Catalogue curation under CC BY 4.0; definitions and English exposition added locally.
\url{https://huggingface.co/datasets/ulamai/UnsolvedMath}.
\item[\hypertarget{source:20002524:2}{S2}] American Institute of Mathematics, $L^2$ harmonic forms in geometry and string theory; notes by Anda Degeratu and Mark Haskins, March 2004, 18 March 2004, Cherkis area-spectrum questions.
\url{https://aimath.org/WWN/l2harmonic/l2harmonic.pdf}.
\end{itemize}
\label{end:20002524}
\end{document}
